\documentclass[a4paper,11pt]{article}

\usepackage[T1]{fontenc}
\usepackage[utf8]{inputenc}
\usepackage{lmodern}
\usepackage{amsmath,amssymb,bm}
\usepackage{booktabs}
\usepackage{array}
\usepackage{geometry}
\usepackage{xcolor}
\usepackage{hyperref}
\usepackage{siunitx}
\usepackage{tikz}
\usepackage{pgfplots}
\pgfplotsset{compat=1.18}

\geometry{margin=2.5cm}

\hypersetup{
  colorlinks=true,
  linkcolor=blue!60!black,
  urlcolor=blue!60!black
}

\title{\textbf{Dy11 Transformer Secondary Voltages\\
Under Single-Phase Voltage Sag on the Primary Side}}
\author{Davide}
\date{\today}

% -----------------------------------------------------------------
\begin{document}
\maketitle
\tableofcontents
\bigskip

% =================================================================
\section{Problem Statement}
\label{sec:problem}

Consider a balanced three-phase primary supply with nominal line-to-neutral
voltage magnitude $V_0 = 1\,\mathrm{pu}$.  A \emph{type~B} single-phase
voltage sag occurs on phase~V: the voltage drops to $0.25\,\mathrm{pu}$,
while phases~U and~W remain at $1\,\mathrm{pu}$.  The transformer is a
\textbf{Dy11} unit with turns ratio $k = N_1/N_2$.

The goal is to determine the secondary line-to-line voltages (and, by
extension, the line-to-neutral voltages of the equivalent star on the
secondary side).

% =================================================================
\section{Primary Phasor Definition}
\label{sec:primary}

Taking the pre-fault phasors as reference (subscript~1 for primary):

\begin{equation}
\underline{V}_{U1} = V_0\,e^{j\,0}
               = V_0,
\qquad
\underline{V}_{V1} = V_0\,e^{-j\,\frac{2\pi}{3}},
\qquad
\underline{V}_{W1} = V_0\,e^{+j\,\frac{2\pi}{3}}.
\end{equation}

After the sag on phase~V with depth $p = 0.25$:

\begin{align}
\underline{V}_{U1} &= V_0, \label{eq:Vu1}\\
\underline{V}_{V1}^{\,\prime} &= p\,V_0\,e^{-j\,\frac{2\pi}{3}}
   = 0.25\,V_0\,e^{-j\,120°}, \label{eq:Vv1}\\
\underline{V}_{W1} &= V_0\,e^{+j\,\frac{2\pi}{3}}
   = V_0\,e^{+j\,120°}. \label{eq:Vw1}
\end{align}

\noindent Using per-unit values with $V_0 = 1$, and expressing in
rectangular form:

\begin{align}
\underline{V}_{U1} &= 1 + j\,0, \\
\underline{V}_{V1}^{\,\prime} &= 0.25\!\left(-\tfrac{1}{2} - j\tfrac{\sqrt{3}}{2}\right)
   = -0.125 - j\,0.2165, \\
\underline{V}_{W1} &= -\tfrac{1}{2} + j\tfrac{\sqrt{3}}{2}
   = -0.5 + j\,0.8660.
\end{align}

% =================================================================
\section{Primary Line-to-Line Voltages}
\label{sec:primary-ll}

The delta winding of a Dy transformer is connected across the
\emph{primary line-to-line} voltages.  With the sag on~V:

\begin{align}
\underline{V}_{UV1} &= \underline{V}_{U1} - \underline{V}_{V1}^{\,\prime}
   = (1+0.125) + j\,0.2165
   = 1.125 + j\,0.2165, \\[4pt]
\underline{V}_{VW1} &= \underline{V}_{V1}^{\,\prime} - \underline{V}_{W1}
   = (-0.125+0.5) + j(-0.2165-0.8660)
   = 0.375 - j\,1.0825, \\[4pt]
\underline{V}_{WU1} &= \underline{V}_{W1} - \underline{V}_{U1}
   = -1.5 + j\,0.8660.
\end{align}

Their magnitudes are:

\begin{align}
|\underline{V}_{UV1}| &= \sqrt{1.125^2 + 0.2165^2}
   = \sqrt{1.2656 + 0.0469} = \sqrt{1.3125} \approx 1.1456\,\mathrm{pu},\\[4pt]
|\underline{V}_{VW1}| &= \sqrt{0.375^2 + 1.0825^2}
   = \sqrt{0.1406 + 1.1718} = \sqrt{1.3124} \approx 1.1456\,\mathrm{pu},\\[4pt]
|\underline{V}_{WU1}| &= \sqrt{1.5^2 + 0.8660^2}
   = \sqrt{2.25 + 0.75} = \sqrt{3} \approx 1.7321\,\mathrm{pu}.
\end{align}

\noindent \textbf{Remark:} Under a balanced supply, all three line-to-line
voltages would equal $\sqrt{3}\,\mathrm{pu} \approx 1.7321\,\mathrm{pu}$.
The sag reduces $V_{UV}$ and $V_{VW}$, while $V_{WU}$ (which does not involve
phase~V directly from $W$ to~$U$) retains its nominal magnitude.

% =================================================================
\section{Dy11 Transformation and Secondary Voltages}
\label{sec:secondary}

\subsection{Dy11 Phase Displacement Convention}

In the \textbf{Dy11} group, the secondary star (y) leads the primary
delta (D) by $11 \times 30° = 330°$, which is equivalent to a
\textbf{lead of $-30°$} (i.e., $+30°$ when measured as the secondary
leading the primary).  The standard convention states:

\begin{equation}
\underline{V}_{a2} = \frac{1}{k}\,\underline{V}_{UV1}\,e^{+j\,30°},
\qquad
\underline{V}_{b2} = \frac{1}{k}\,\underline{V}_{VW1}\,e^{+j\,30°},
\qquad
\underline{V}_{c2} = \frac{1}{k}\,\underline{V}_{WU1}\,e^{+j\,30°},
\label{eq:dy11}
\end{equation}

where $k$ is the primary-to-secondary turns ratio.  Setting $k=1$ for
per-unit analysis (the secondary base voltage is chosen so that rated
secondary equals 1\,pu when the primary is balanced at 1\,pu):

\begin{equation}
\frac{1}{k} = \frac{1}{\sqrt{3}}
\end{equation}

because the primary line-to-line base equals $\sqrt{3}$ times the
primary phase base.  The $+30°$ rotation is the defining
characteristic of the Dy11 group.

\subsection{Secondary Line-to-Neutral Phasors}

Applying \eqref{eq:dy11} with the $+30°$ rotation operator
$e^{+j30°} = \frac{\sqrt{3}}{2} + j\frac{1}{2}$:

\paragraph{Phase a (mapped from UV):}
\begin{align}
\underline{V}_{a2} &= \frac{1}{\sqrt{3}}\,(1.125 + j\,0.2165)\,
   \left(\frac{\sqrt{3}}{2} + j\,\frac{1}{2}\right) \notag\\
&= \frac{1}{\sqrt{3}}\,
   \bigl[(1.125\cdot\tfrac{\sqrt{3}}{2} - 0.2165\cdot\tfrac{1}{2})
   + j\,(1.125\cdot\tfrac{1}{2} + 0.2165\cdot\tfrac{\sqrt{3}}{2})\bigr] \notag\\
&= \frac{1}{\sqrt{3}}\,
   \bigl[(0.9743 - 0.1083) + j\,(0.5625 + 0.1875)\bigr] \notag\\
&= \frac{1}{\sqrt{3}}\,(0.8660 + j\,0.75) \notag\\
&\approx 0.5 + j\,0.4330.
\end{align}

Magnitude: $|\underline{V}_{a2}| = \sqrt{0.5^2 + 0.4330^2}
= \sqrt{0.25 + 0.1875} = \sqrt{0.4375} \approx \mathbf{0.6614\,pu}$.

\paragraph{Phase b (mapped from VW):}
\begin{align}
\underline{V}_{b2} &= \frac{1}{\sqrt{3}}\,(0.375 - j\,1.0825)\,
   \left(\frac{\sqrt{3}}{2} + j\,\frac{1}{2}\right) \notag\\
&= \frac{1}{\sqrt{3}}\,
   \bigl[(0.375\cdot\tfrac{\sqrt{3}}{2} + 1.0825\cdot\tfrac{1}{2})
   + j\,(0.375\cdot\tfrac{1}{2} - 1.0825\cdot\tfrac{\sqrt{3}}{2})\bigr] \notag\\
&= \frac{1}{\sqrt{3}}\,
   \bigl[(0.3248 + 0.5413) + j\,(0.1875 - 0.9375)\bigr] \notag\\
&= \frac{1}{\sqrt{3}}\,(0.8661 - j\,0.75) \notag\\
&\approx 0.5 - j\,0.4330.
\end{align}

Magnitude: $|\underline{V}_{b2}| = \sqrt{0.5^2 + 0.4330^2} \approx \mathbf{0.6614\,pu}$.

\paragraph{Phase c (mapped from WU):}
\begin{align}
\underline{V}_{c2} &= \frac{1}{\sqrt{3}}\,(-1.5 + j\,0.8660)\,
   \left(\frac{\sqrt{3}}{2} + j\,\frac{1}{2}\right) \notag\\
&= \frac{1}{\sqrt{3}}\,
   \bigl[(-1.5\cdot\tfrac{\sqrt{3}}{2} - 0.8660\cdot\tfrac{1}{2})
   + j\,(-1.5\cdot\tfrac{1}{2} + 0.8660\cdot\tfrac{\sqrt{3}}{2})\bigr] \notag\\
&= \frac{1}{\sqrt{3}}\,
   \bigl[(-1.2990 - 0.4330) + j\,(-0.75 + 0.75)\bigr] \notag\\
&= \frac{1}{\sqrt{3}}\,(-1.7320 + j\,0) \notag\\
&= -1.0 + j\,0.
\end{align}

Magnitude: $|\underline{V}_{c2}| = \mathbf{1.0\,pu}$.

% =================================================================
\section{Summary of Results}
\label{sec:summary}

Table~\ref{tab:results} collects all primary and secondary voltages in per-unit.

\begin{table}[h!]
\centering
\caption{Primary and secondary voltages under type~B sag on phase~V
         ($p = 0.25\,\mathrm{pu}$), Dy11 transformer, $k{=}1\,\mathrm{pu}$.}
\label{tab:results}
\begin{tabular}{lcccc}
\toprule
\textbf{Quantity} & \textbf{Phasor (pu)} & \textbf{Magnitude (pu)}
                  & \textbf{Angle (deg)} & \textbf{Nominal (pu)}\\
\midrule
$V_{U1}$  & $1{+}j0$               & 1.000 & $0°$    & 1.000 \\
$V_{V1}'$ & $-0.125{-}j0.2165$     & 0.250 & $-120°$ & 1.000 \\
$V_{W1}$  & $-0.5{+}j0.8660$       & 1.000 & $+120°$ & 1.000 \\
\midrule
$V_{UV1}$ & $1.125{+}j0.2165$      & 1.146 & $+10.9°$  & $\sqrt{3}$\\
$V_{VW1}$ & $0.375{-}j1.0825$      & 1.146 & $-70.9°$  & $\sqrt{3}$\\
$V_{WU1}$ & $-1.5{+}j0.8660$       & 1.732 & $+150°$   & $\sqrt{3}$\\
\midrule
$V_{a2}$  & $0.5{+}j0.4330$        & 0.661 & $+40.9°$  & 1.000 \\
$V_{b2}$  & $0.5{-}j0.4330$        & 0.661 & $-40.9°$  & 1.000 \\
$V_{c2}$  & $-1.0{+}j0$            & 1.000 & $+180°$   & 1.000 \\
\bottomrule
\end{tabular}
\end{table}

\subsection{Physical Interpretation}

\begin{itemize}
  \item \textbf{Phase c on the secondary is unaffected} ($|V_{c2}|=1\,\mathrm{pu}$)
        because it is derived from the $WU$ primary line-to-line voltage,
        which involves only phases U and W — both at rated voltage.

  \item \textbf{Phases a and b are equally reduced} to $\approx 0.661\,\mathrm{pu}$
        (i.e., $\tfrac{1}{\sqrt{3}} \cdot \sqrt{1.3125} \approx 0.661$),
        since both $V_{UV}$ and $V_{VW}$ are affected symmetrically by the
        sag on V.

  \item The \textbf{angular spread} of the secondary phasors is no longer
        $120°$: phases a and b are only $81.8°$ apart instead of $120°$,
        while the angle from b to c and from c to a are both $\approx 139.1°$.
        This unbalanced phase portrait is the characteristic signature of a
        Dy-transformed type~B sag.

  \item In terms of \emph{symmetrical components}, the primary unbalanced
        set contains a negative-sequence component introduced by the sag.
        The delta primary winding acts as a \emph{zero-sequence trap} (no
        zero-sequence current circulates in the secondary star), but the
        negative-sequence component passes through undistorted (with the
        $30°$ rotation of the Dy11 group).  Consequently the secondary
        contains both positive- and negative-sequence components, but
        \emph{no zero-sequence component}.
\end{itemize}

% =================================================================
\section{Symmetrical-Component Verification}
\label{sec:symcomp}

The secondary unbalanced set can be decomposed as
$\underline{V}_{abc} = \underline{V}^{+} + \underline{V}^{-}$
(no zero sequence due to the delta primary).
Applying the Fortescue transformation:

\begin{equation}
\begin{pmatrix}V^0\\V^+\\V^-\end{pmatrix}
= \frac{1}{3}
\begin{pmatrix}1&1&1\\1&\alpha&\alpha^2\\1&\alpha^2&\alpha\end{pmatrix}
\begin{pmatrix}V_{a2}\\V_{b2}\\V_{c2}\end{pmatrix},
\qquad \alpha = e^{j\,120°}.
\end{equation}

With $V_{a2} = 0.5+j0.4330$, $V_{b2} = 0.5-j0.4330$,
$V_{c2} = -1$:

\begin{align}
V^0 &= \tfrac{1}{3}(V_{a2}+V_{b2}+V_{c2})
     = \tfrac{1}{3}(0.5+0.5-1) = 0 \quad\checkmark \\[4pt]
V^+ &= \tfrac{1}{3}(V_{a2}+\alpha\,V_{b2}+\alpha^2 V_{c2}).
\end{align}

Recalling $\alpha = -\tfrac{1}{2}+j\tfrac{\sqrt{3}}{2}$ and
$\alpha^2 = -\tfrac{1}{2}-j\tfrac{\sqrt{3}}{2}$:

\begin{align}
\alpha\,V_{b2} &= \left(-\tfrac{1}{2}+j\tfrac{\sqrt{3}}{2}\right)
                 (0.5-j0.4330)
               = (-0.25+0.3750)+j(-0.2165+0.4330\cdot\tfrac{\sqrt{3}}{2}\cdot 2)\notag\\
               &= (-0.25+0.3750)+j(-0.2165+0.3750) \notag\\
               &\approx 0.125 + j\,0.6495, \notag\\
\alpha^2 V_{c2} &= \left(-\tfrac{1}{2}-j\tfrac{\sqrt{3}}{2}\right)(-1)
                = 0.5 + j\,0.8660, \notag\\[4pt]
V^+ &= \tfrac{1}{3}\bigl[(0.5+j0.4330)+(0.125+j0.6495)+(0.5+j0.8660)\bigr] \notag\\
    &= \tfrac{1}{3}(1.125 + j\,1.9485)
     \approx 0.375 + j\,0.6495
     \approx 0.75\,e^{j\,60°}, \\[4pt]
|V^+| &\approx 0.750\,\mathrm{pu}.
\end{align}

Similarly, $|V^-| \approx 0.250\,\mathrm{pu}$ (the negative-sequence
magnitude equals the per-unit sag depth residual).  This confirms the
well-known result for a type~B sag transformed through a Dy connection:

\begin{equation}
\boxed{
|V^+| = \frac{2+p}{3} = \frac{2+0.25}{3} = 0.75\,\mathrm{pu},
\qquad
|V^-| = \frac{1-p}{3} = \frac{1-0.25}{3} = 0.25\,\mathrm{pu}.
}
\end{equation}

% =================================================================
\section{Conclusions}
\label{sec:conclusions}

For a \textbf{type~B single-phase voltage sag} ($V_V = 0.25\,\mathrm{pu}$,
$V_U = V_W = 1\,\mathrm{pu}$) on the primary side of a \textbf{Dy11}
transformer, the secondary voltages are:

\begin{center}
\begin{tabular}{lcc}
\toprule
\textbf{Secondary phase} & \textbf{Magnitude (pu)} & \textbf{Note} \\
\midrule
$a$ & 0.661 & Reduced (involves sagged phase V via UV) \\
$b$ & 0.661 & Reduced (involves sagged phase V via VW) \\
$c$ & 1.000 & Unaffected (WU path; V not involved) \\
\bottomrule
\end{tabular}
\end{center}

\noindent The delta primary \emph{converts} the single-phase primary sag
(type~B) into a \textbf{two-phase secondary sag} (type~C/D family),
with no zero-sequence propagation.  The negative-sequence component at
the secondary equals $|V^-|=(1-p)/3 = 0.25\,\mathrm{pu}$, which drives
negative-sequence currents in any downstream motor or power-electronic
load — a critical input for ride-through and control design.

% =================================================================
\section{Inverse Problem: Primary Sag Required to Obtain a Type~B Sag at the Secondary}
\label{sec:inverse}

\subsection{Problem Statement}

We now seek the \emph{primary} phase voltages such that the \textbf{secondary}
presents a type~B sag:

\begin{equation}
V_{U2} = 1\,\mathrm{pu},\qquad
V_{V2} = 0.25\,\mathrm{pu}\;(\angle{-120°}),\qquad
V_{W2} = 1\,\mathrm{pu}\;(\angle{+120°}).
\label{eq:sec_target}
\end{equation}

\subsection{Secondary Phasors in Rectangular Form}

\begin{align}
\underline{V}_{U2} &= 1 + j\,0, \\
\underline{V}_{V2} &= 0.25\,e^{-j120°} = -0.125 - j\,0.2165, \\
\underline{V}_{W2} &= 1\,e^{+j120°}   = -0.5   + j\,0.8660.
\end{align}

\subsection{Secondary Line-to-Line Voltages}

\begin{align}
\underline{V}_{UV2} &= \underline{V}_{U2} - \underline{V}_{V2}
   = 1.125 + j\,0.2165, \quad |\underline{V}_{UV2}| \approx 1.146\,\mathrm{pu},\\
\underline{V}_{VW2} &= \underline{V}_{V2} - \underline{V}_{W2}
   = 0.375 - j\,1.0825, \quad |\underline{V}_{VW2}| \approx 1.146\,\mathrm{pu},\\
\underline{V}_{WU2} &= \underline{V}_{W2} - \underline{V}_{U2}
   = -1.5   + j\,0.8660, \quad |\underline{V}_{WU2}| = \sqrt{3}\,\mathrm{pu}.
\end{align}

\subsection{Inverse Dy11 Mapping}

The Dy11 group applies a $+30°$ rotation and a $\sqrt{3}$~scaling from
primary line-to-line to secondary line-to-neutral (Section~\ref{sec:secondary}).
The inverse mapping is therefore a $-30°$ rotation and a $\sqrt{3}$~upscaling:

\begin{equation}
\underline{V}_{UV1} = \sqrt{3}\;\underline{V}_{U2}\,e^{-j\,30°},
\qquad
\underline{V}_{VW1} = \sqrt{3}\;\underline{V}_{V2}\,e^{-j\,30°},
\qquad
\underline{V}_{WU1} = \sqrt{3}\;\underline{V}_{W2}\,e^{-j\,30°}.
\label{eq:inv_dy11}
\end{equation}

The $-30°$ rotation operator is
$e^{-j30°} = \tfrac{\sqrt{3}}{2} - j\tfrac{1}{2}$.

\paragraph{Primary line UV:}
\begin{align}
\underline{V}_{UV1}
  &= \sqrt{3}\,(1+j\,0)\!\left(\tfrac{\sqrt{3}}{2}-j\tfrac{1}{2}\right)
   = \sqrt{3}\!\left(\tfrac{\sqrt{3}}{2}-j\tfrac{1}{2}\right)
   = \tfrac{3}{2} - j\tfrac{\sqrt{3}}{2}
   \approx 1.5 - j\,0.8660.
\end{align}
Magnitude: $|\underline{V}_{UV1}| = \sqrt{1.5^2+0.8660^2}=\sqrt{3}\,\mathrm{pu}$. \quad
Angle: $-30°$.

\paragraph{Primary line VW:}
\begin{align}
\underline{V}_{VW1}
  &= \sqrt{3}\,(-0.125-j\,0.2165)\!\left(\tfrac{\sqrt{3}}{2}-j\tfrac{1}{2}\right)\notag\\
  &= \sqrt{3}\!\left[(-0.125\cdot\tfrac{\sqrt{3}}{2}-0.2165\cdot\tfrac{1}{2})
     +j\,(-0.125\cdot(-\tfrac{1}{2})+(-0.2165)\cdot\tfrac{\sqrt{3}}{2})\right]\notag\\
  &= \sqrt{3}\!\left[(-0.1083-0.1083)+j\,(0.0625-0.1875)\right]\notag\\
  &= \sqrt{3}\,(-0.2165-j\,0.125)\notag\\
  &\approx -0.375 - j\,0.2165.
\end{align}
Magnitude: $\sqrt{3}\cdot 0.25 \approx 0.433\,\mathrm{pu}$. \quad
Angle: $-150°$.

\paragraph{Primary line WU:}
\begin{align}
\underline{V}_{WU1}
  &= \sqrt{3}\,(-0.5+j\,0.8660)\!\left(\tfrac{\sqrt{3}}{2}-j\tfrac{1}{2}\right)\notag\\
  &= \sqrt{3}\!\left[(-0.5\cdot\tfrac{\sqrt{3}}{2}-0.8660\cdot(-\tfrac{1}{2}))
     +j\,((-0.5)\cdot(-\tfrac{1}{2})+0.8660\cdot\tfrac{\sqrt{3}}{2})\right]\notag\\
  &= \sqrt{3}\!\left[(-0.4330+0.4330)+j\,(0.25+0.75)\right]\notag\\
  &= \sqrt{3}\,(0+j\,1.0)\notag\\
  &\approx 0 + j\,1.7321.
\end{align}
Magnitude: $\sqrt{3}\cdot 1.0 \approx 1.732\,\mathrm{pu}$. \quad
Angle: $+90°$.

\subsection{Primary Phase-to-Neutral Voltages}

The primary line-to-line voltages are related to the primary
phase-to-neutral voltages by:

\begin{equation}
\underline{V}_{UV1} = \underline{V}_{U1} - \underline{V}_{V1},\quad
\underline{V}_{VW1} = \underline{V}_{V1} - \underline{V}_{W1},\quad
\underline{V}_{WU1} = \underline{V}_{W1} - \underline{V}_{U1}.
\end{equation}

This is a $3\times 3$ singular system (only two equations are independent
for a three-phase system without neutral connection).  Imposing the
conventional constraint $\underline{V}_{U1}+\underline{V}_{V1}+\underline{V}_{W1}=0$
and solving:

\begin{align}
\underline{V}_{U1} &= \tfrac{1}{3}\!\left(\underline{V}_{UV1}-\underline{V}_{WU1}\right),\\
\underline{V}_{V1} &= \tfrac{1}{3}\!\left(\underline{V}_{VW1}-\underline{V}_{UV1}\right),\\
\underline{V}_{W1} &= \tfrac{1}{3}\!\left(\underline{V}_{WU1}-\underline{V}_{VW1}\right).
\end{align}

Substituting the computed line-to-line phasors:

\paragraph{Phase U:}
\begin{align}
\underline{V}_{U1}
  &= \tfrac{1}{3}\bigl[(1.5-j\,0.8660)-(0+j\,1.7321)\bigr]\notag\\
  &= \tfrac{1}{3}(1.5 - j\,2.5981)\notag\\
  &\approx 0.5 - j\,0.8660
   = e^{-j\,60°}.
\end{align}
$|\underline{V}_{U1}|=1\,\mathrm{pu}$,\quad $\angle\underline{V}_{U1}=-60°$.

\paragraph{Phase V:}
\begin{align}
\underline{V}_{V1}
  &= \tfrac{1}{3}\bigl[(-0.375-j\,0.2165)-(1.5-j\,0.8660)\bigr]\notag\\
  &= \tfrac{1}{3}(-1.875+j\,0.6495)\notag\\
  &\approx -0.625 + j\,0.2165.
\end{align}
$|\underline{V}_{V1}|=\sqrt{0.625^2+0.2165^2}=\sqrt{0.3906+0.0469}=\sqrt{0.4375}\approx
\mathbf{0.661\,\mathrm{pu}}$,\quad
$\angle\underline{V}_{V1}=180°-\arctan(0.2165/0.625)\approx 180°-19.1°=+160.9°$.

\paragraph{Phase W:}
\begin{align}
\underline{V}_{W1}
  &= \tfrac{1}{3}\bigl[(0+j\,1.7321)-(-0.375-j\,0.2165)\bigr]\notag\\
  &= \tfrac{1}{3}(0.375+j\,1.9486)\notag\\
  &\approx 0.125 + j\,0.6495.
\end{align}
$|\underline{V}_{W1}|=\sqrt{0.125^2+0.6495^2}=\sqrt{0.0156+0.4219}=\sqrt{0.4375}\approx
\mathbf{0.661\,\mathrm{pu}}$,\quad
$\angle\underline{V}_{W1}=\arctan(0.6495/0.125)\approx+79.1°$.

\subsection{Summary: Required Primary Sag}

\begin{table}[h!]
\centering
\caption{Primary voltages required to produce a type~B sag
         ($V_{V2}{=}0.25\,\mathrm{pu}$) at the secondary of a Dy11
         transformer.}
\label{tab:inv_results}
\begin{tabular}{lcccl}
\toprule
\textbf{Primary phase} & \textbf{Phasor (pu)} & \textbf{Magnitude (pu)}
   & \textbf{Angle (deg)} & \textbf{Comment} \\
\midrule
U & $0.5 - j\,0.866$ & 1.000 & $-60°$  & Nominal magnitude, shifted $-60°$ \\
V & $-0.625 + j\,0.2165$ & 0.661 & $+160.9°$ & \textbf{Reduced to $0.661\,\mathrm{pu}$, shifted} \\
W & $0.125 + j\,0.6495$ & 0.661 & $+79.1°$  & \textbf{Reduced to $0.661\,\mathrm{pu}$, shifted} \\
\bottomrule
\end{tabular}
\end{table}

\subsection{Physical Interpretation of the Inverse Result}

Several key observations follow from this inverse analysis:

\begin{itemize}
  \item The required primary sag is a \textbf{type~C sag}: \emph{two phases}
        are reduced in magnitude (V and W, both to $\approx 0.661\,\mathrm{pu}$),
        while one phase (U) retains nominal magnitude.  This is exactly the
        dual of the forward transformation: just as a Dy connection turns a
        type~B primary sag into a type~C/D secondary sag, it turns a
        type~C primary sag into a type~B secondary sag.

  \item The entire primary phasor set is \textbf{rotated by $-60°$} with
        respect to the standard reference position.  This is a consequence
        of the $-30°$ inverse Dy11 mapping applied to a secondary
        reference frame that itself started from $0°$.  In practice, the
        phase angles depend on the chosen reference frame; what matters
        physically is only the \emph{relative} magnitudes and angles
        between the three phases.

  \item The \textbf{sag depth at the primary} corresponding to the
        $0.25\,\mathrm{pu}$ secondary type~B sag is:
        \begin{equation}
          p_{\mathrm{prim}} = 0.661\,\mathrm{pu}
          = \sqrt{\frac{1+p_{\mathrm{sec}}^2 + p_{\mathrm{sec}}}{3}}
          = \sqrt{\frac{1+0.25^2+0.25}{3}} = \sqrt{\frac{1.3125}{3}}
          \approx 0.661\,\mathrm{pu},
        \end{equation}
        confirming the analytical formula for Dy-transformed type~B/C sag pairs.

  \item The \textbf{symmetrical-component content} of the required primary
        type~C sag is $|V^+_1|\approx0.7737\,\mathrm{pu}$ and
        $|V^-_1|\approx0.1130\,\mathrm{pu}$.  These values differ from
        those at the secondary (Section~\ref{sec:symcomp}: $|V^+_2|=0.75$,
        $|V^-_2|=0.25$~pu) because the Dy transformation does \emph{not}
        preserve sequence magnitudes when the sag \emph{type} changes
        (type~C primary $\leftrightarrow$ type~B secondary); it only
        preserves magnitudes for the \emph{same} sag type propagating
        unchanged through the transformer.
\end{itemize}

\noindent In summary: to observe a type~B sag ($U{=}1$, $V{=}0.25$,
$W{=}1$~pu) at the \emph{secondary} of a Dy11 transformer, one must apply
a \textbf{type~C sag with depth $\approx 0.661\,\mathrm{pu}$ on two phases}
(V and W) at the primary, with the phasor frame rotated by $-60°$ relative
to the secondary reference.

% =================================================================
\section{$D$--$Q$ Decomposition of the Primary Type~C Sag}
\label{sec:dq}

\subsection{Definition of the $D$--$Q$ Frame}

In the synchronous $dq$ reference frame rotating at angular frequency
$\omega$, a three-phase voltage set is decomposed into:
\begin{itemize}
  \item a \textbf{positive-sequence} phasor $\underline{V}^+= D_p + jQ_p$,
        co-rotating with the frame;
  \item a \textbf{negative-sequence} phasor $\underline{V}^-= D_n + jQ_n$,
        counter-rotating at $2\omega$ in the frame.
\end{itemize}
The Fortescue decomposition gives:
\begin{equation}
\underline{V}^+ = \frac{1}{3}\bigl(V_U + \alpha\,V_V + \alpha^2 V_W\bigr),
\qquad
\underline{V}^- = \frac{1}{3}\bigl(V_U + \alpha^2 V_V + \alpha\,V_W\bigr),
\qquad \alpha = e^{j\,120°}.
\end{equation}

\subsection{Canonical Type~C Frame ($V_U = 1\angle 0°$)}

The type~C sag with $p = 0.661\,\mathrm{pu}$ and $V_U$ as phase reference:
\begin{equation}
\underline{V}_U = 1,\qquad
\underline{V}_V = 0.661\,e^{-j120°},\qquad
\underline{V}_W = 0.661\,e^{+j120°}.
\end{equation}

For a type~C sag with two phases symmetrically reduced to $p$, the
closed-form Fortescue result is:
\begin{equation}
\underline{V}^+ = \frac{1+2p}{3},\qquad
\underline{V}^- = \frac{1-p}{3},\qquad
\underline{V}^0 = 0.
\label{eq:typeC_symcomp}
\end{equation}

With $p = 0.661$:
\begin{align}
\underline{V}^+ &= \frac{1+2\times0.661}{3} = \frac{2.322}{3}
  = 0.7740\,\mathrm{pu}\;\angle\,0°, \\[4pt]
\underline{V}^- &= \frac{1-0.661}{3} = \frac{0.339}{3}
  = 0.1130\,\mathrm{pu}\;\angle\,0°.
\end{align}

Both phasors are \emph{real} in the canonical frame (no imaginary part),
so:
\begin{equation}
\boxed{
D_p = 0.7740\,\mathrm{pu},\quad Q_p = 0,\quad
D_n = 0.1130\,\mathrm{pu},\quad Q_n = 0.
}
\label{eq:dq_canonical}
\end{equation}

\noindent \textbf{Remark}: the absence of $Q$ components in the canonical
frame is the defining property of a type~C sag — the sag is purely real
when $V_U$ is the $D$-axis reference.  Any non-zero $Q$ component signals
either a non-zero fault impedance angle, a measurement frame offset, or a
different sag type.

\subsection{Rotated Frame ($V_U = 1\angle{-60°}$, as required at the Dy11 primary)}

When the primary phasor set is expressed in the rotated frame of
Section~\ref{sec:inverse} (where $V_{U1} = e^{-j60°}$), the
$dq$ reference axis itself is rotated by $\theta_0 = -60°$ with respect
to the canonical axis.  The sequence phasors transform as:
\begin{equation}
\underline{V}^+_{\theta_0} = \underline{V}^+\,e^{j\theta_0}
  = 0.7740\,e^{-j60°},\qquad
\underline{V}^-_{\theta_0} = \underline{V}^-\,e^{-j\theta_0}
  = 0.1130\,e^{+j60°}.
\end{equation}

\noindent (Positive sequence rotates with the frame; negative sequence
rotates opposite.)  In rectangular form:
\begin{align}
\underline{V}^+_{\theta_0}
  &= 0.7740\!\left(\tfrac{1}{2} - j\tfrac{\sqrt{3}}{2}\right)
   = 0.3870 - j\,0.6703, \\[4pt]
\underline{V}^-_{\theta_0}
  &= 0.1130\!\left(\tfrac{1}{2} + j\tfrac{\sqrt{3}}{2}\right)
   = 0.0565 + j\,0.0978.
\end{align}

Thus, in the rotated primary frame:
\begin{equation}
\boxed{
D_p = 0.3870\,\mathrm{pu},\quad Q_p = -0.6703\,\mathrm{pu},\quad
D_n = 0.0565\,\mathrm{pu},\quad Q_n = +0.0978\,\mathrm{pu}.
}
\label{eq:dq_rotated}
\end{equation}

\subsection{Summary Table}

\begin{table}[h!]
\centering
\caption{Symmetrical components and $D$--$Q$ decomposition of the
         primary type~C sag ($p = 0.661\,\mathrm{pu}$) in both reference frames.}
\label{tab:dq}
\begin{tabular}{lcccccc}
\toprule
\textbf{Frame} & $|V^+|$ & $|V^-|$ & $D_p$ & $Q_p$ & $D_n$ & $Q_n$ \\
               & (pu)    & (pu)    & (pu)  & (pu)  & (pu)  & (pu)  \\
\midrule
Canonical ($V_U = 1\angle 0°$)         & 0.7740 & 0.1130 & $+0.7740$ & $0$       & $+0.1130$ & $0$       \\
Rotated ($V_U = 1\angle{-60°}$, Dy11)  & 0.7740 & 0.1130 & $+0.3870$ & $-0.6703$ & $+0.0565$ & $+0.0978$ \\
\bottomrule
\end{tabular}
\end{table}

\subsection{Relevance for Control Design}

In a vector-controlled drive or grid-tied converter operating under this
type~C primary sag, the $D$--$Q$ components have direct implications:

\begin{itemize}
  \item $D_p$ and $Q_p$ feed the \emph{positive-sequence current
        controller} (inner $dq$ loop).  A drop in $D_p$ from 1.0 to
        $0.774\,\mathrm{pu}$ reduces the available active power headroom.

  \item $D_n$ and $Q_n$ generate a \emph{double-frequency ripple} at
        $2\omega$ in the positive-sequence $dq$ frame.  Without a
        dedicated negative-sequence controller (e.g.\ a second $dq$ loop
        at $-\omega$, or a notch filter at $2\omega$ on the current
        references), this ripple appears as a $100\,\mathrm{Hz}$ (or
        $120\,\mathrm{Hz}$) disturbance on the $d$- and $q$-axis currents.

  \item In the canonical frame, $Q_n = 0$ means the negative-sequence
        voltage has no reactive component along the $q$-axis: the ripple
        is \emph{in-phase} with the $D$-axis and affects primarily the
        active-power channel.  In the rotated frame both channels are
        disturbed.

  \item The ratio $|V^-|/|V^+| = 0.1130/0.7740 \approx 14.6\%$
        is the \emph{voltage unbalance factor} (VUF) at the primary
        terminals.  IEC~61000-2-2 sets a planning level of $2\%$ for LV
        networks and $1\%$ for HV; a VUF of $14.6\%$ is therefore a
        severe unbalance condition requiring ride-through or protection
        action.
\end{itemize}

\end{document}

